% META: {"id":"1NSI-TC-CO-022","chapitre":"1NSI-TYPES-CONSTRUITS","type_objet":"corrige","capacites":["1NSI-TYPES-CONSTRUITS-C2"],"mode_creation":"ex_nihilo","sources_inspiration":[],"status":"approved","exercice_ref":"1NSI-TC-EX-022","origin":"needs_review","status_history":["needs_review","audited","accepted","approved"],"release_acceptance":"RELEASE_OWNER_BATCH_ACCEPTANCE","acceptance_closure_digest":"sha256:9b3ccf9a81c5520fbb7e03b7b2d3e7bf2057a02834b6d908bf3be5f49f3b553f","acceptance_receipt_digest":"sha256:4fad86f0282e0fa4e0419cca1ad6379ae7af5a280c04a4a90880f90a1c044242"}
% BEGIN-VERIFY
% eleves = ["Alice", "Bob", "Clara"]
% notes = [14, 8, 17]
% resultats = []
% for i in range(len(eleves)):
%     if notes[i] >= 10:
%         resultats.append((eleves[i], notes[i]))
% assert resultats == [("Alice", 14), ("Clara", 17)]
% END-VERIFY
\begin{corrige}{1NSI-TC-EX-022}
\begin{enumerate}
  \item Les affichages produits sont :
  \begin{itemize}
    \item \lstinline{Alice : 14} (car $14 \geq 10$).
    \item \lstinline{Clara : 17} (car $17 \geq 10$).
  \end{itemize}
  Bob n'apparaît pas car sa note ($8$) est inférieure à 10.

  \item Les deux tableaux sont des \textbf{tableaux parallèles} : l'élément d'indice \lstinline{i} dans \lstinline{eleves} correspond à l'élément de même indice dans \lstinline{notes}. Par exemple, \lstinline{eleves[0]} (\lstinline{"Alice"}) a pour note \lstinline{notes[0]} (\lstinline{14}).

  \item On utilise \lstinline{range(len(eleves))} car on a besoin de l'indice \lstinline{i} pour accéder simultanément à \lstinline{eleves[i]} et \lstinline{notes[i]}. Avec \lstinline{for e in eleves}, on obtiendrait les noms mais on ne pourrait pas accéder à la note correspondante dans \lstinline{notes}.
\end{enumerate}

Vérification : les résultats ont été confirmés par exécution.
\end{corrige}
